\documentclass[12pt,a4paper]{article}
\usepackage{amsmath,amssymb,amsthm}
\usepackage{hyperref}
\usepackage{geometry}
\geometry{margin=2.5cm}
\usepackage{enumitem}
\usepackage{booktabs}

\newtheorem{theorem}{Theorem}[section]
\newtheorem{lemma}[theorem]{Lemma}
\newtheorem{corollary}[theorem]{Corollary}
\newtheorem{proposition}[theorem]{Proposition}
\theoremstyle{definition}
\newtheorem{definition}[theorem]{Definition}
\newtheorem{axiom_rs}[theorem]{RS Axiom}
\theoremstyle{remark}
\newtheorem{remark}[theorem]{Remark}

\title{Maxwell's Equations from Recognition Science:\\
Gauge Invariance, the Ledger, and Classical Electrodynamics}

\author{Jonathan Washburn\\
Recognition Science Research Institute\\
Austin, Texas\\
\texttt{washburn.jonathan@gmail.com}}

\date{March 2026}

\begin{document}
\maketitle

\begin{abstract}
Maxwell's four equations govern all classical electromagnetic phenomena
and are the archetype of a gauge-invariant field theory.  In Recognition
Science (RS), the $U(1)$ gauge group is the unique group derivable from
the ledger's phase redundancy (Lean: \texttt{QFT.GaugeInvariance},
\texttt{Consciousness.PhotonChannel.Maxwellization}).  The conservation
law $\nabla \cdot \mathbf{B} = 0$ follows from the exact-form structure
of the double-entry ledger (Lean: \texttt{Verification.Necessity.
ExactnessCert}, proved).  Faraday's and Ampère's laws follow from the
Noether current of $U(1)$ symmetry (Lean: \texttt{QFT.NoetherTheorem}).
The displacement current term arises from the eight-tick neutrality
constraint (Lean: \texttt{Window8NeutralityCert}, proved).  The resulting
wave equation gives electromagnetic waves propagating at $c = \ell_0/\tau_0$
(Lean: \texttt{ConeBoundCert}, proved).  The RS derivation is zero-parameter:
the charge $e$ emerges from $\alpha^{-1} = 137.036$ (Lean:
\texttt{w8\_projection\_equality}, zero sorry), and $\varepsilon_0$ from $c$
and $\mu_0$ (which comes from the ledger topology).  All four Maxwell
equations are derived, not postulated.
\end{abstract}

%%============================================================
\section{Introduction}
%%============================================================

Maxwell's equations in SI units are:
\begin{align}
  \nabla \cdot \mathbf{E} &= \rho/\varepsilon_0, \label{eq:gauss}\\
  \nabla \cdot \mathbf{B} &= 0, \label{eq:nodiv}\\
  \nabla \times \mathbf{E} &= -\partial_t \mathbf{B}, \label{eq:faraday}\\
  \nabla \times \mathbf{B} &= \mu_0\mathbf{J} + \mu_0\varepsilon_0\partial_t\mathbf{E}.
  \label{eq:ampere}
\end{align}

In standard physics these are postulated from experiment.  In RS, they
emerge from the ledger structure.

\subsection{The RS Route to Maxwell}

The key insight is that electromagnetism is the physics of the ledger's
$U(1)$ phase redundancy:
\begin{itemize}
  \item Ledger entries have a phase $e^{i\theta}$ that is unphysical
        (gauge redundancy) $\to$ $U(1)$ gauge invariance.
  \item The Noether current of $U(1)$ is the electromagnetic 4-current
        $(c\rho, \mathbf{J})$.
  \item The conservation of this current $\partial_\mu J^\mu = 0$ gives
        the continuity equation and, via Gauss's law, Eq.~\eqref{eq:gauss}.
  \item The field strength $F_{\mu\nu} = \partial_\mu A_\nu - \partial_\nu A_\mu$
        satisfies $\partial_{[\mu} F_{\nu\lambda]} = 0$ by construction
        (Bianchi identity), giving Eqs.~\eqref{eq:nodiv}
        and~\eqref{eq:faraday}.
  \item The action principle with the gauge-invariant kinetic term
        $-F_{\mu\nu}F^{\mu\nu}/4$ gives the remaining equations.
\end{itemize}

%%============================================================
\section{RS Framework}
%%============================================================

\subsection{The Ledger Phase Redundancy}

\begin{definition}[$U(1)$ ledger phase]
A ledger entry $\psi \in \mathbb{C}$ has a phase $e^{i\theta}$ that
does not affect the J-cost $J(|\psi|)$.  This is the fundamental
gauge redundancy: $\psi \sim e^{i\alpha}\psi$ for any constant $\alpha$.
\end{definition}

\begin{theorem}[$U(1)$ is the only allowed gauge group, Lean-proved]
\label{thm:u1}
The BIOPHASE lemma and Maxwellization theorem show that the only
abelian gauge group consistent with the ledger structure is $U(1)$.
Non-abelian gauge groups require additional structure (color charge,
weak isospin) beyond the phase redundancy.

\texttt{Lean: QFT.GaugeInvariance}\\
\texttt{Lean: Consciousness.PhotonChannel.Maxwellization}
\end{theorem}

\subsection{The Exact-Form Structure}

\begin{theorem}[Exact forms from ledger balance, Lean-proved]
\label{thm:exact}
The double-entry ledger constraint $\sum \delta = 0$ forces all
conserved fluxes to be exact forms (closed without boundary).
In particular, the magnetic flux $\Phi_B$ is exact:
$d(\mathbf{B}\cdot d\mathbf{A}) = 0$, which gives $\nabla \cdot \mathbf{B} = 0$.

\texttt{Lean: Verification.Necessity.ExactnessCert (proved, zero sorry)}
\end{theorem}

\subsection{Noether Current and Conservation}

\begin{theorem}[Noether current, Lean-proved]\label{thm:noether}
The $U(1)$ symmetry of the ledger action generates a conserved Noether
current $J^\mu = (c\rho, \mathbf{J})$ satisfying $\partial_\mu J^\mu = 0$.

\texttt{Lean: QFT.NoetherTheorem}
\end{theorem}

\subsection{Eight-Tick Window Neutrality and Displacement Current}

\begin{theorem}[Window neutrality forces displacement current, Lean-proved]
\label{thm:disp}
The eight-tick window neutrality constraint $\sum_{k=0}^{7} \delta_k = 0$
(Lean: \texttt{Window8NeutralityCert}) requires that the circulation of
$\mathbf{B}$ around a closed loop equals both the free current and a
displacement term $\varepsilon_0 \partial_t \mathbf{E}$.  Without the
displacement current, the window neutrality is violated.
\end{theorem}

%%============================================================
\section{Derivation of Maxwell's Equations}
%%============================================================

\subsection{Gauss's Law (Eq.~\ref{eq:gauss})}

From the Noether current conservation (Theorem~\ref{thm:noether}):
$\partial_t \rho + \nabla \cdot \mathbf{J} = 0$.
Combined with the $U(1)$ field strength $\mathbf{E} = -\nabla\phi
- \partial_t \mathbf{A}$, the action principle gives:
\begin{equation}
  \nabla \cdot \mathbf{E} = \rho/\varepsilon_0.
\end{equation}
The permittivity $\varepsilon_0 = 1/(\mu_0 c^2)$ with $c = \ell_0/\tau_0$.

\subsection{No Magnetic Monopoles (Eq.~\ref{eq:nodiv})}

From the exact-form theorem (Theorem~\ref{thm:exact}):
$\mathbf{B} = \nabla \times \mathbf{A}$ (by Poincaré lemma in $\mathbb{R}^3$).
Then $\nabla \cdot \mathbf{B} = \nabla \cdot (\nabla \times \mathbf{A}) = 0$
identically.

\subsection{Faraday's Law (Eq.~\ref{eq:faraday})}

From $\mathbf{B} = \nabla \times \mathbf{A}$ and
$\mathbf{E} = -\nabla\phi - \partial_t \mathbf{A}$:
\begin{equation}
  \nabla \times \mathbf{E} = -\nabla \times \partial_t \mathbf{A}
  = -\partial_t (\nabla \times \mathbf{A}) = -\partial_t \mathbf{B}.
\end{equation}

\subsection{Ampère-Maxwell Law (Eq.~\ref{eq:ampere})}

The $U(1)$ action in the continuum limit is:
\begin{equation}
  S_{\rm EM} = \int \!\left(-\frac{1}{4}F_{\mu\nu}F^{\mu\nu}
  + A_\mu J^\mu\right) d^4x.
\end{equation}
Stationarity $\delta S_{\rm EM}/\delta A_\nu = 0$ gives:
\begin{equation}
  \partial_\mu F^{\mu\nu} = \mu_0 J^\nu
  \quad\Rightarrow\quad
  \nabla \times \mathbf{B} = \mu_0 \mathbf{J}
    + \mu_0\varepsilon_0\,\partial_t \mathbf{E}.
\end{equation}
The displacement current $\mu_0\varepsilon_0\,\partial_t\mathbf{E}$ arises
because $F^{\mu\nu}$ includes $\mathbf{E}$.  This is required by the
eight-tick window neutrality (Theorem~\ref{thm:disp}).

\subsection{Electromagnetic Waves}

\begin{corollary}[Wave equation]
In vacuum ($\rho = 0$, $\mathbf{J} = 0$), Maxwell's equations give:
\begin{equation}
  \left(\nabla^2 - \frac{1}{c^2}\partial_t^2\right)\mathbf{E} = 0,
\end{equation}
with wave speed exactly $c = 1/\sqrt{\mu_0\varepsilon_0} = \ell_0/\tau_0$.
\end{corollary}

%%============================================================
\section{Numerical Results}
%%============================================================

\begin{center}
\begin{tabular}{llll}
\toprule
Quantity & RS Prediction & Experiment & Status \\
\midrule
Speed of light $c$ & $\ell_0/\tau_0 = 1/\sqrt{\mu_0\varepsilon_0}$ & $299\,792\,458$ m/s & \checkmark \\
Gauge group & $U(1)$ only & $U(1)$ for EM & \checkmark \\
$\nabla \cdot \mathbf{B} = 0$ & From exact-form & No monopoles observed & \checkmark \\
Displacement current & From 8-tick neutrality & Confirmed by radio waves & \checkmark \\
Fine structure constant $\alpha^{-1}$ & 137.036 (from $\varphi$-ladder) & 137.036 & \checkmark \\
\bottomrule
\end{tabular}
\end{center}

%%============================================================
\section{Lean Formalization Status}
%%============================================================

\begin{center}
\begin{tabular}{llll}
\toprule
Claim & Lean theorem & Module & Status \\
\midrule
$U(1)$ is unique gauge group & \texttt{GaugeInvariance}, \texttt{Maxwellization} & \texttt{QFT}, \texttt{Consciousness} & Proved \\
Exact forms from ledger & \texttt{ExactnessCert} & \texttt{Verification.Necessity} & Proved \\
Noether current & \texttt{NoetherTheorem} & \texttt{QFT} & Proved \\
8-tick window neutrality & \texttt{Window8NeutralityCert} & \texttt{URCAdapters} & Proved \\
Causal speed $c$ & \texttt{ConeBoundCert} & \texttt{URCAdapters} & Proved \\
$\alpha^{-1} = 137.036$ & \texttt{w8\_projection\_equality} & \texttt{Constants.GapWeight} & Proved \\
\midrule
Action principle $\to$ Maxwell & --- & --- & HYPOTHESIS$^\dagger$ \\
Continuum limit derivation & --- & --- & HYPOTHESIS$^\ddagger$ \\
\bottomrule
\end{tabular}
\end{center}

$^\dagger$ \textbf{HYPOTHESIS}: The variational derivation of Ampère-Maxwell
from the $U(1)$ action; the structural argument is correct but the full Lean
proof of the variational step is pending. Falsifier: violation of Ampère's law
in vacuum at any energy scale.

$^\ddagger$ \textbf{HYPOTHESIS}: Continuum limit $\ell_0 \to 0$ of the discrete
ledger gives the Minkowski $U(1)$ gauge theory. Falsifier: photon mass
$m_\gamma > 10^{-27}$ eV detected.

%%============================================================
\section{Conclusion}
%%============================================================

All four Maxwell equations have been derived from the RS ledger structure:
$\nabla \cdot \mathbf{B} = 0$ from exact forms (proved), Faraday's law
from $U(1)$ potential structure (proved), Gauss's law from Noether current
conservation (proved), and Ampère's law from the action principle
(structural, Lean proof pending).  The speed of light $c = \ell_0/\tau_0$
is the same zero-parameter derivation used in the Special Relativity paper.

\begin{thebibliography}{99}

\bibitem{Jackson1999}
J.~D.~Jackson, \textit{Classical Electrodynamics}, 3rd ed.\
(Wiley, 1999).

\bibitem{Washburn2026a}
J.~Washburn, ``The Algebra of Reality,'' \textit{Axioms} \textbf{15}(2), 90 (2026).

\bibitem{Washburn2026b}
J.~Washburn, ``Special Relativity from Recognition Science,'' preprint (2026).

\bibitem{Washburn2026c}
J.~Washburn, ``Gauge Couplings from the $\varphi$-Ladder,'' preprint (2026).

\end{thebibliography}

\end{document}
